This dissertation presents a layer-wise study of predictive-coding (PC) regimes
relative to backpropagation (BP) under comparable computational conditions. The work keeps
four analytical dimensions separate: similarity of final behavior, similarity
of internal mechanism, computational cost, and the possibility of replacing the
remaining canonical iterative suffix before normal completion with a registered
bounded analytic completion while preserving the required response. To connect
internal state observations to a compute-control decision, two bounded theoretical
constructs are introduced. PC-TREF defines task-relative state equivalence, full
sufficiency of a diagnostic representation on the registered domain, and a
narrower selective sufficiency condition for a specific action. PC-CATM describes
correction formation, transport, and cancellation channels. QWake-PC operationalizes
these foundations as a compute-control architecture; only its bounded FixedPred
instantiation, QWake-FP, is evaluated empirically. In QWake-FP, the early action
selects the registered analytic completion
\nolinkurl{fixedpred_eta1_wavefront_completion_v1} instead of the remaining
canonical iterative sweeps; the complete exact suffix
\nolinkurl{complete_suffix_stage2_baseline_v1} remains the reference and fallback
path.

The experimental program is organized as a sequence of preregistered and versioned
stages. Stage~1 and Stage~2 compare BP, Exact, FixedPred, and Strict under controlled
conditions. Stage~3A examines layer-wise gradients and representations, while
Stage~3B localizes state-inference cost and evaluates exact computational-replacement
candidates through registered equivalence gates. The final QWake-FP branch separates
three questions: whether the sealed trajectory surface contains preterminal records
whose states make the registered analytic-completion action task-admissible; whether
a non-zero subset can be selectively recognized from the pre-action state with zero
observed dangerous accepts on the frozen calibration surface; and whether that
decision yields positive aggregate net saving under frozen full decision-cost
accounting.

Within the investigated conditions, FixedPred showed higher observed similarity to
BP than Strict in gradient direction and learned representations, while its
early-layer gradient norm was substantially reduced. Stage~3B profiling associated
a substantial share of the measured PC cost with state inference; subsequent B1/B2
candidates passed their registered equivalence gates, and the matched profiling
matrix was completed without retries. The registered engineering-continuation
screen then qualified 0 of 16 configurations in each of the four
B1/B2-by-FixedPred/Strict groups, and both candidates received
\texttt{reject\_or\_revise}; this bounded negative result is not interpreted as a
superiority claim for the baseline.

In QWake-FP, the sealed surface contained preterminal records whose states made the
registered analytic-completion action task-admissible, and non-zero selective
pre-action recognizability with zero observed dangerous accepts was found within the
frozen family of 2,625 scalar rules on 756 calibration records. Of these, 264 rules
had non-zero coverage with zero dangerous accepts, and the highest-coverage rule
among rules with zero observed dangerous accepts covered 216 of 756 calibration
records (28.57\%). That rule uses only \texttt{compute\_step >= 5}; therefore C08
establishes a temporal boundary equivalent to a fixed prefix on this surface, not
input-dependent adaptivity or superiority of PC-CATM features. Structural
reconciliation shows that these 216 records comprise 108 preterminal records at
step 5 with one remaining sweep and 108 records at the terminal boundary step 6.
Thus, 28.57\% remains the registered full-surface coverage and is not interpreted
as preterminal analytic-action coverage. No rule achieved positive aggregate net
saving under the frozen full decision-cost accounting. The full cost assigned to
the highest-coverage zero-danger rule was overwhelmingly observer-dominated; no C2
rule freeze was established, and confirmatory C3 therefore remained closed.

The negative economic result is specific to the frozen full-cost accounting used in
C2 and does not measure the marginal runtime cost of a minimal recognizer. The
dissertation therefore supports the presence of preterminal records whose states
make the registered analytic-completion action task-admissible and non-zero selective
pre-action recognizability with zero observed dangerous accepts on the frozen
calibration surface; rejects the existence of a rule with zero observed dangerous accepts,
non-zero coverage, and positive aggregate net saving within the frozen C2 family under
that accounting regime; and leaves marginal implementation viability as a separate
untested estimand.

\noindent\textbf{Keywords:} predictive coding, backpropagation, layer-wise analysis,
computational cost, analytic completion, traceability.
